% 総目録（紙見本・背景パターン編）．
% 31 種の背景パターンを系統ごとに並べ，濃度・色味・寸法の効き方を示す．
% 扉の帯に敷いた実寸の見え方は examples/magazine/surikumi-patterns.tex にある．
% 本見本は「どれがどれか」を一覧で引くためのものである．
\documentclass[lualatex,paper=b5j,fontsize=10pt]{jlreq}

\usepackage[twocolumn=false,footer=false]{surikumi}

% --- 見本用の局所部品 -------------------------------------------------
\newcommand{\CatName}[1]{{\ttfamily\fontsize{7.4}{9}\selectfont #1}}

\newcommand{\CatHead}[2]{%
  \par\addvspace{1.1\baselineskip}%
  \noindent{\sffamily\gtfamily\bfseries\fontsize{10.5}{12.5}\selectfont #1}%
  \hspace{.7em}{\color{SMrule}\fontsize{7.6}{9}\selectfont #2}\par
  \nobreak\vspace{.8mm}%
  \noindent{\color{SMink}\rule{\linewidth}{.5pt}}\par
  \nobreak\addvspace{.5\baselineskip}%
}

\newcommand{\CatUse}[1]{%
  \par\addvspace{.2\baselineskip}%
  \noindent{\color{SMrule}\fontsize{7.6}{10}\selectfont #1}\par
  \nobreak\addvspace{.35\baselineskip}%
}

% パターン1件．見本と鍵の名前を縦に積む．\MMPatternSwatch は
% \MathMagazineSetup で選んだパターンと濃度をそのまま使うので，
% 直前に pattern を差し替えてから呼ぶ．
% \MMPatternSwatch の幅は TikZ の座標へそのまま入るので，
% \dimexpr を渡さず，あらかじめ長さへ入れて確定させておく．
\newlength{\CatCell}
\setlength{\CatCell}{30mm}
\newlength{\CatSwatchW}
\setlength{\CatSwatchW}{27mm}

\newcommand{\CatSwatch}[1]{%
  \begin{minipage}[t]{\CatCell}
    \MathMagazineSetup{pattern=#1}%
    \MMPatternSwatch[13mm]{\CatSwatchW}\par
    \vspace{.8mm}
    \CatName{#1}\par
    \vspace{2.6mm}
  \end{minipage}\hspace{2mm}%
}

% 濃度・色味・寸法の比較用．鍵をそのまま見出しにする．
% 4 枚を 1 行に並べるので，本文幅 154mm に収まる寸法にしておく．
\newlength{\CatTunedW}
\setlength{\CatTunedW}{34mm}
\newcommand{\CatTuned}[2]{%
  \begin{minipage}[t]{\CatTunedW}
    \MathMagazineSetup{#1}%
    \MMPatternSwatch[13mm]{\CatTunedW}\par
    \vspace{.8mm}
    \CatName{#2}\par
    \vspace{2.6mm}
  \end{minipage}\hspace{1.5mm}%
}

\MathMagazineSetup{
  feature={},
  series={総目録},
  series-position=left,
  pattern=asanoha,
  title={背景パターン31種},
  author={数理組版},
  author-font=mincho,
  author-position=right,
  title-fill=black,
  deck={記事扉のバナーに敷く背景パターンを，系統ごとに一覧する．いずれも公有の意匠を pgf のパスで独自に作図したもので，参考紙面の図案の複製ではない．}
}

\begin{document}

\MMTitle

上の帯が \CatName{pattern=asanoha} を実寸で敷いた例である．
以下の見本はすべて \CatName{\string\MMPatternSwatch} による縮小見本で，
\CatName{pattern-scale=0.34} を掛けて図柄が1単位分収まるようにしてある．
実寸の見え方は \CatName{examples/magazine/surikumi-patterns.tex} で確かめる．

\MathMagazineSetup{pattern-scale=0.34}

% =====================================================================
\CatHead{幾何}{12種}

\CatUse{当初からある系統．タイル寸法は宣言時に固定されており，
pattern-scale は効かない．}

\noindent
\CatSwatch{scales}\CatSwatch{triangles}\CatSwatch{diamonds}\CatSwatch{chevrons}\par
\noindent
\CatSwatch{weave}\CatSwatch{hexagons}\CatSwatch{pinwheels}\CatSwatch{zigzags}\par
\noindent
\CatSwatch{checker}\CatSwatch{parquet}\CatSwatch{origami}\CatSwatch{ribbons}\par

\CatUse{既定は scales である．}

% =====================================================================
\CatHead{和柄}{10種}

\CatUse{inherently coloured pattern として宣言してあり，タイル内部の階調を
そのまま出力へ残す．したがって pattern-accent と pattern-tone のうち
インク色の指定は効かず，白ベールだけが効く．pattern-scale も効かない．}

\noindent
\CatSwatch{asanoha}\CatSwatch{seigaiha}\CatSwatch{shippou}\CatSwatch{kikkou}\par
\noindent
\CatSwatch{sayagata}\CatSwatch{kagome}\CatSwatch{yagasuri}\CatSwatch{tatewaku}\par
\noindent
\CatSwatch{mitsukuzushi}\CatSwatch{ajiro}\par

\CatUse{線画の6種（asanoha, seigaiha, shippou, kikkou, sayagata, kagome）が薄く，
面塗りの4種（yagasuri, tatewaku, mitsukuzushi, ajiro）が濃い．
白抜きの題字を載せるときは薄い側を選ぶ．}

% =====================================================================
\CatHead{数学的モチーフ}{9種}

\CatUse{新しい鍵駆動の宣言を使っており，この9種にだけ pattern-scale が効く．}

\noindent
\CatSwatch{truchet}\CatSwatch{isocubes}\CatSwatch{sierpinski}\CatSwatch{lattice}\par
\noindent
\CatSwatch{moire}\CatSwatch{envelopes}\CatSwatch{girih}\CatSwatch{packing}\par
\noindent
\CatSwatch{spirals}\par

\CatUse{薄いものから濃いものへ並べると，
lattice ＜ sierpinski ＜ envelopes ＜ truchet ＜ spirals ＜ packing ＜
girih ＜ moire ＜ isocubes の順になる．}

% =====================================================================
\CatHead{濃度}{pattern-tone}

\CatUse{normal が従来の見え方である．題字が読みにくいときだけ pale か light へ落とす．
strong は白ベールを外すので，題字を載せない帯にだけ使う．}

\noindent
\CatTuned{pattern=isocubes,pattern-tone=pale}{pattern-tone=pale}%
\CatTuned{pattern=isocubes,pattern-tone=light}{pattern-tone=light}%
\CatTuned{pattern=isocubes,pattern-tone=normal}{normal（既定）}%
\CatTuned{pattern=isocubes,pattern-tone=strong}{pattern-tone=strong}\par

% =====================================================================
\CatHead{色味}{pattern-accent}

\CatUse{背景の色味だけに使う．本文の色には使わない．
濃度を落とした状態と組み合わせるのが基本である．}

\noindent
\CatTuned{pattern=girih,pattern-tone=light,pattern-accent=SMcoverBlue}%
  {accent=SMcoverBlue}%
\CatTuned{pattern=sierpinski,pattern-tone=pale,pattern-accent=SMcoverOrange}%
  {accent=SMcoverOrange}%
\CatTuned{pattern=truchet,pattern-tone=light,pattern-accent=SMcoverViolet}%
  {accent=SMcoverViolet}%
\CatTuned{pattern=lattice,pattern-tone=normal,pattern-accent=SMcoverCyan}%
  {accent=SMcoverCyan}\par

% =====================================================================
\CatHead{寸法}{pattern-scale}

\CatUse{数学的モチーフの9種にだけ効く．
幾何12種と和柄10種へ指定しても黙って無視される（誤りではない）．}

\noindent
\CatTuned{pattern=moire,pattern-tone=normal,pattern-accent=none,pattern-scale=0.24}%
  {moire, scale=0.24}%
\CatTuned{pattern=moire,pattern-scale=0.5}{moire, scale=0.5}%
\CatTuned{pattern=packing,pattern-scale=0.24}{packing, scale=0.24}%
\CatTuned{pattern=packing,pattern-scale=0.5}{packing, scale=0.5}\par

% =====================================================================
\CatHead{バナー以外の部品}{\string\MMPatternRule}

\CatUse{節の区切りに使う細い帯である．高さは省略可能引数で変えられる
（既定 3.2mm）．見出しの代わりには使わない．}

\MathMagazineSetup{pattern=sayagata,pattern-tone=light,pattern-scale=1}

\MMPatternRule

\CatUse{既定の高さ．}

\MMPatternRule[6mm]

\CatUse{\string\MMPatternRule[6mm]．}

行の中へ置く見本は \CatName{\string\MMPatternSwatch} である．
\MathMagazineSetup{pattern-scale=0.34}%
truchet \MathMagazineSetup{pattern=truchet}\MMPatternSwatch{11mm}，
girih \MathMagazineSetup{pattern=girih}\MMPatternSwatch{11mm}，
spirals \MathMagazineSetup{pattern=spirals}\MMPatternSwatch{11mm}
のように使う．短縮名は \CatName{\string\MMPRule} と
\CatName{\string\MMSwatch} である．

\end{document}
